\section{Intrinsic uniqueness of the original all-degree norm factor}

\begin{lemma}\label{lem:native-tc-actual-positive-degree-section}
\lean{EnvelopeProperty.NativeTestConfiguration.TestConfiguration.exists_positive_gaussPowerSectionNorm}
\leanok
\uses{thm:native-test-configuration-actual-gauss-finite-norm-family}
For original geometrically ample $L$, every actual original Gauss
point has a positive-degree original total-space section with
strictly positive true total-model norm.
\end{lemma}
\begin{proof}
The genuine original ampleness producer gives a threshold $N$
and an actual finite original section family in each degree at
least $N$, chosen before every original point, with strictly
positive maximum true norm at every actual original Gauss
point. Use degree $N+1$, which is positive and at least $N$.
A finite maximum of nonnegative values is positive only if
one of its actual members is positive. Choose that original
member at the specified point. This supplies the claimed
positive-degree section without asserting anything in degree one.
\end{proof}

\begin{theorem}\label{thm:native-tc-unique-positive-all-power-factor}
\lean{EnvelopeProperty.NativeTestConfiguration.TestConfiguration.existsUnique_weightedLaurentPowerSectionNorm_powerFactor}
\leanok
\uses{lem:native-tc-actual-positive-degree-section,thm:native-tc-coherent-all-power-section-norms}
For the original ample $L$ there is exactly one positive real
factor $c$ satisfying the true original norm identity
$\lVert s\rVert_{\mathcal X,r,x}=c^n W_{n,x,r}(s)$
for every nonnegative degree and every original section.
The existing factor is the native valuation of an actual frame
transition unit; uniqueness does not redefine it by degree-one
section values.
\end{theorem}
\begin{proof}
The coherent original norm theorem constructs such a positive
factor from an actual native unit before all degrees and
sections. Suppose $c$ and $d$ both satisfy the asserted
identities. Choose the actual positive-degree section of the
preceding lemma, in degree $m>0$. Its true norm is positive.
Its weighted original coefficient maximum is therefore
nonzero: if it were zero, the identity for $c$ would force
the true norm to be zero. The two identities give
$c^mW=d^mW$. Cancel the nonzero $W$ to obtain $c^m=d^m$.
The positive-power map is injective on nonnegative real
numbers for $m>0$, so $c=d$. This proves uniqueness.
\end{proof}

\begin{corollary}\label{cor:native-tc-original-frame-factor-independence}
\lean{EnvelopeProperty.NativeTestConfiguration.TestConfiguration.gaussCompatibleNormFactor_eq_of_frames}
\leanok
\uses{thm:native-tc-unique-positive-all-power-factor,def:native-tc-coherent-positive-norm-factor,thm:native-tc-coherent-all-power-section-norms}
The actual native positive factors constructed using any two
original opens and original frames around the same proper
center of $x$ are equal, for the original ample $L$.
\end{corollary}
\begin{proof}
Each actual original center-open frame defines a positive
native transition valuation. The coherent all-degree theorem
shows that each satisfies the same original all-degree,
all-section norm identity. The preceding uniqueness theorem
identifies them. Original total-model and original coefficient
norms remain the same in both identities, so no frame matching
condition is needed.
\end{proof}
